\section{신호 처리 원시 연산자의 수학적 계약}
\label{sec:operators}
원시 연산자는 수치 함수와 함께 입력 스키마, 출력 스키마, 상태 불변식, 시간 지지,
초기화 규칙과 품질 전파 규칙을 제공한다. 본 절은 전기생리ㆍ광학 신호의 핵심 연산을
이 계약의 형태로 정의한다.

\subsection{단위 변환과 재기준화}
단위 변환은 채널별 선형 또는 affine 사상으로 표현한다. 전극 전압의 재기준화는
가중치 $w\in\mathbb R^C$, $w^\top\mathbf 1=1$을 사용하여 다음과 같이 계산한다.
\begin{equation}
 A_w=I-\mathbf 1w^\top,\qquad Y=A_wX.
 \label{eq:reref}
\end{equation}
연산자는 출력 기준을 $w$와 입력 채널 식별자의 조합으로 기록한다.
$A_w\mathbf 1=0$이므로 공통 성분에 해당하는 영공간을 형성한다.
후속 공분산 연산은 이 선형 변환과 자신의 정규화 계약을 함께 사용한다.
채널 선택과 순서 변경은 명시적인 선택 행렬 $P$로 표현하며 $Y=PX$와 출력 스키마 $P\mathcal C$를 생성한다.

\subsection{필터와 표본률 변환}
길이 $L$인 FIR 필터를 다음 상태 전이로 정의한다.
\begin{equation}
 y_c[n]=\sum_{\ell=0}^{L-1}h[\ell]x_c[n-\ell].
 \label{eq:fir}
\end{equation}
상태는 채널별 이전 $L-1$개 표본이다. 초기 상태의 값과 초기 유효성 구간을 계약에 기록한다.
계수 $h$가 무차원이면 출력은 입력 전압의 단위를 유지한다.
출력의 대표 시각은 $t_n$으로 정하고, 시간 지지는 $t_{n-L+1}$부터 $t_n$까지의 입력 구간으로 구성한다.
선형 위상 계수가 $h[\ell]=h[L-1-\ell]$을 만족하면 군지연은 $(L-1)/(2f_s)$이다.
실행기는 이 신호 지연과 계산 완료 시각을 각각 기록한다.

IIR 필터는 상태공간 표현을 사용한다.
\begin{equation}
 z_{n+1}=Az_n+Bx_n,\qquad y_n=Cz_n+Dx_n.
\end{equation}
필터 프로필은 계수의 수치 형식, 실현 방식, 초기 상태, 포화 정책과 안정성 조건을 가진다.
$\rho(A)<1$인 선형 시불변 실현은 유계 입력에 대한 상태의 유계성을 제공한다.
실행기는 선택한 계수와 부동소수점 실현을 버전으로 고정한다.

유리수 비율 $p/q$의 재표본화는 입력 격자와 출력 격자를 연결한다.
\begin{equation}
 f_{\mathrm{out}}=\frac pq f_{\mathrm{in}},\qquad
 y[m]=\sum_n x[n]h[mq-np],\qquad p,q\in\mathbb N_{>0}.
 \label{eq:resample}
\end{equation}
계약은 재구성 필터 $h$, 위상 원점, 경계 초기값, anti-alias 대역을 지정한다.
출력 계보는 각 표본에 기여한 입력 구간과 재표본화 설정을 연결한다.
정렬 연산자는 대상 시계를 선택하고, 재표본화 연산자는 대상 격자를 선택한다.

\subsection{창과 사건 기준 구간}
끝점 기준 창은 시작 위상 $n_0$, 길이 $N$, hop $H$, $N,H\in\mathbb N_{>0}$를 가진다.
\begin{equation}
 X_k=[x[n_0+kH],x[n_0+kH+1],\ldots,x[n_0+kH+N-1]]
      \in\mathbb R^{C\times N}.
 \label{eq:window}
\end{equation}
출력 값 형식은 $\operatorname{Epoch}[D,C,N]$이며,
출력 대표 시각은 $t_{n_0+kH+N-1}$이고, 연속 출력의 명목 간격은 $H/f_s$이다.
표본 첫 시각과 끝 시각의 차이는 $(N-1)/f_s$이다.
원시 구간의 표본 집합과 품질 마스크가 창의 유효성을 결정한다.
사건 기준 구간은 사건 $e$의 시각 $t_e$와 상대 구간 $[a,b)$를 사용하여
$[t_e+a,t_e+b)$의 표본을 선택한다. 입력 진행 경계가 구간의 끝에 도달하면 창을 완결한다.
중심 창은 대표 시각과 전방 관측 길이를 연산자 효과에 포함한다.

\subsection{전력 스펙트럼과 무차원 특징}
실수 입력과 길이 $N\ge1$의 창 함수 $w[n]$를 사용한다.
창 함수는 $\sum_{n=0}^{N-1}w[n]^2>0$을 만족한다.
이산 푸리에 변환을
\begin{equation}
 \widehat X_c[k]=\sum_{n=0}^{N-1}w[n]x_c[n]
                 e^{-2\pi\mathrm{i}kn/N}
\end{equation}
로 정의한다. 한쪽 스펙트럼의 정규화는
\begin{equation}
 P_c[k]=\frac{\eta_k|\widehat X_c[k]|^2}
                  {f_s\sum_{n=0}^{N-1}w[n]^2},\qquad
 f_k=\frac{kf_s}{N}
 \label{eq:psd}
\end{equation}
를 $k=0,\ldots,\lfloor N/2\rfloor$에 사용한다.
$\eta_k$는 DC와 짝수 $N$의 Nyquist 성분에서 1, 나머지 양의 주파수에서 2이다.
입력 단위가 $u$이면 $P_c[k]$의 단위는 $u^2/\mathrm{Hz}$이다.
주파수 대역 $B$의 전력과 로그 특징은
\begin{equation}
 E_{c,B}=\sum_{k:f_k\in B}P_c[k]\frac{f_s}{N},\qquad
 \varphi_{c,B}=\log\frac{E_{c,B}+E_0}{E_{\mathrm{ref}}}
\end{equation}
로 정의한다. 양의 상수 $E_0,E_{\mathrm{ref}}$는 $u^2$ 단위를 가지며,
특징값 $\varphi_{c,B}$는 무차원이다. 언어는 대역 경계, 창 함수, 정규화와 기준 전력을
특징 스키마에 기록한다.

\subsection{공분산과 공간 특징}
평균 제거 행렬을 $H_N=I_N-\mathbf 1\mathbf 1^\top/N$으로 두면 표본 공분산은
\begin{equation}
 \widehat\Sigma=\frac{1}{N-1}XH_NX^\top,\qquad N\ge2.
\end{equation}
공분산 연산은 동일한 물리량과 단위를 갖는 채널 묶음을 입력으로 받는다.
일반적인 다중 물리량 입력은 선언한 채널별 정규화로 무차원 벡터를 생성한 후 공분산을 계산한다.
양정치 행렬을 요구하는 후속 연산을 위해 다음 정규화를 정의한다.
\begin{equation}
 \begin{aligned}
 s^2&=\max\left(\frac{\tr\widehat\Sigma}{C},\sigma_0^2\right),
       &\sigma_0^2&>0,\\
 \Sigma_\lambda&=(1-\lambda)\widehat\Sigma+\lambda s^2I_C,
       &0&<\lambda\le1.
 \end{aligned}
 \label{eq:spd}
\end{equation}
$\sigma_0^2$는 공분산과 같은 단위를 갖는다.

\begin{proposition}[정규화 공분산의 양정치성]
$\widehat\Sigma\succeq0$, $\sigma_0^2>0$, $0<\lambda\le1$이면
$\Sigma_\lambda\succ0$이다.
\end{proposition}
\begin{proof}
임의의 $z\ne0$에 대해
\[
 z^\top\Sigma_\lambda z
 =(1-\lambda)z^\top\widehat\Sigma z+\lambda s^2\lVert z\rVert^2
 \ge\lambda\sigma_0^2\lVert z\rVert^2>0.
\]
따라서 양정치 조건이 성립한다.
\end{proof}

CSP 연산자는 두 클래스의 공분산에서 공간 가중치를 학습한다\cite{csp}.
본 언어의 CSP 계약은 정규화된 클래스 공분산 $\Sigma_1,\Sigma_2$와 다음 고유문제를 사용한다.
\begin{equation}
 \Sigma_1w_j=\lambda_j(\Sigma_1+\Sigma_2)w_j,
 \qquad w_j^\top(\Sigma_1+\Sigma_2)w_j=1.
\end{equation}
선택한 고유벡터, 입력 채널 순서와 기준을 모델 상태에 저장한다.
선택한 고유벡터의 유한 인덱스 집합을 $J$, $|J|\ge1$로 둔다.
현재 창의 정규화 공분산으로부터 다음 로그 특징을 계산한다.
\begin{equation}
 e_j=w_j^\top\Sigma_\lambda w_j,\qquad
 \varphi_j=\log\frac{e_j}{\sum_{\ell\in J}e_\ell},\qquad j\in J.
 \label{eq:csp-feature}
\end{equation}
고유벡터의 정규화와 $\Sigma_\lambda\succ0$에서 $e_j>0$을 얻는다.
분산의 비율은 무차원이며 로그 특징의 정의역을 만족한다.

Riemannian 특징은 양정치 공분산의 기하를 이용한다\cite{riemannian}.
학습 집합에서 구한 기준 $R\succ0$에 대해
\begin{equation}
 Z=\log(R^{-1/2}\Sigma_\lambda R^{-1/2}),\qquad
 \varphi=\operatorname{svec}(Z)
 \label{eq:tangent}
\end{equation}
를 사용한다. $\operatorname{svec}$는 대각 원소와 $\sqrt2$배의 상삼각 비대각 원소를
벡터로 구성하여 Frobenius 내적을 보존한다. 식~\eqref{eq:tangent}의 로그 인수는 무차원이다.
기준 $R$과 채널 스키마는 모델 계보에 속한다.

\subsection{광학 측정 변환}
광강도 $I_\lambda(t)>0$와 양의 기준 $I_{0,\lambda}$에서 광밀도 변화를
\begin{equation}
 \Delta\mathrm{OD}_\lambda(t)
 =-\log\frac{I_\lambda(t)}{I_{0,\lambda}}
\end{equation}
로 정의한다. 광경로와 흡광 계수를 포함한 행렬 $A$를 사용하면
\begin{equation}
 \Delta\mathrm{OD}=A\,\Delta c+\xi,
 \qquad
 \widehat{\Delta c}=(A^\top WA)^{-1}A^\top W\Delta\mathrm{OD}
 \label{eq:mbll}
\end{equation}
를 얻는다. 이 계약은 $W\succ0$, $\rank(A)$가 추정 농도 성분 수와 같다는 조건을 사용한다.
광학 연산자는 선택한 파장, source--detector 거리, 차등 광경로 인자,
흡광 계수의 단위와 로그 밑을 $A$의 메타데이터로 기록한다.
이 구성은 fNIRS 분석의 광밀도ㆍ농도 변환 관계를 따른다\cite{homer}.
SNIRF의 채널 대응 정보와 NIRx의 원시ㆍ파생 출력은 동일한 광학 스키마로 연결한다\cite{snirf,nirx_formats}.

\subsection{수치 실행 계약}
각 연산자는 자료형, 반올림 정책, 유효 입력 영역, 출력 오차 기준을 선언한다.
참조 연산 $F$와 실행 커널 $\widetilde F$의 비교에는
\begin{equation}
 \lVert\widetilde F(x)-F(x)\rVert
 \le\epsilon_{\mathrm{abs}}+\epsilon_{\mathrm{rel}}\lVert F(x)\rVert
 \label{eq:numeric}
\end{equation}
형태의 조건을 사용한다. 정수 디코더는 패킷 규격에 대응하는 정확한 값 관계를 사용하고,
부동소수점 커널은 선택한 입력 영역과 자료형에 맞는 수치 조건을 사용한다.
전이 상태의 유효성 검사와 품질 정책은 각 연산자의 실행 결과에 함께 기록한다.
